\documentclass{article}
\usepackage{amsmath}
\title{Eurobot 2026 Robot PCB Documentation}
\author{Game Team SLSC}
\date{\today}

\begin{document}
\maketitle

\section{Introduction}
This document contains the documentation for the Eurobot 2026 Robot PCB.

We selected 0805 package resistors as minimal case as they represent the minimum comfortable size for hand soldering. \\
Pulldown/up resistors will not be considered in the power dissipation calculations as their values are greater than $4.7k\Omega$, resulting in negligible power dissipation.

\section{Batt Input}
\section{Batt Input Analysis}

\subsection{MOSFET (Q2: IRF4905)}
The IRF4905 acts as a high-side P-channel switch for reverse polarity protection. Its key specifications are:
\begin{itemize}
    \item Max continuous drain current: $I_{D} = -74\text{A}$ (at $T_C = 25^\circ\text{C}$)
    \item Max power dissipation: $P_{D} = 150\text{W}$ (at $T_C = 25^\circ\text{C}$)
    \item $R_{DS(on)}$ (static): $0.02\Omega$ (at $V_{GS} = -10\text{V}$, $I_{D} = -19\text{A}$)
    \item Max Gate-Source Voltage: $V_{GS(max)} = \pm 20\text{V}$
\end{itemize}

The maximum current drawn by the robot is $I_{max} = 15\text{A}$. 
At $25^\circ\text{C}$ junction temperature, the power dissipation is:
$$P = R_{DS(on)} \cdot I^2 = 0.02\Omega \cdot (15\text{A})^2 = 4.5\text{W}$$

\textbf{Thermal Correction:} In real-world conditions, $R_{DS(on)}$ increases with temperature. At $T_J = 125^\circ\text{C}$, $R_{DS(on)}$ is approximately $1.6 \times$ its nominal value ($\approx 0.032\Omega$), leading to $P \approx 7.2\text{W}$. 
With a $R_{\theta JA}$ of $62^\circ\text{C/W}$ for a TO-220 package without a heatsink, the temperature rise would exceed $400^\circ\text{C}$. Therefore, a \textbf{heatsink} or significant copper pouring on the PCB is mandatory.

\subsection{Zener Diode (D2: 1N4744A)}
The 1N4744A is used to protect the \textbf{Gate-Source junction} of the MOSFET, not the downstream load. 
\begin{itemize}
    \item Zener voltage: $V_{Z} = 15\text{V}$ (at $I_{Z} = 18\text{mA}$)
    \item Max power dissipation: $P_{Z} = 1\text{W}$
    \item Current limit: $I_{Zmax} = \frac{P_{Z}}{V_{Z}} = \frac{1\text{W}}{15\text{V}} = 66.7\text{mA}$
\end{itemize}

In this circuit, the current through the Zener is limited by $R2$ ($10\text{k}\Omega$). If the battery voltage is $24\text{V}$:
$$I_{Z} = \frac{V_{BATT} - V_{Z}}{R2} = \frac{24\text{V} - 15\text{V}}{10,000\Omega} = 0.9\text{mA}$$
This is well within the $66.7\text{mA}$ limit. The Zener ensures $V_{GS}$ never exceeds $15\text{V}$, staying safely below the $20\text{V}$ limit of the IRF4905.

\subsection{TVS Diode (D3: SMBJ18A)}
The SMBJ18A provides transient voltage suppression for the \textbf{entire system}.
\begin{itemize}
    \item Reverse Standoff Voltage: $V_{RWM} = 18\text{V}$ (The diode does not conduct at normal $12\text{V}$ battery levels).
    \item Peak Pulse Power: $P_{PP} = 600\text{W}$ (for $10/1000\mu\text{s}$ pulses).
\end{itemize}
This component is critical for absorbing inductive spikes from motors or "hot-plugging" the battery, preventing the destruction of downstream regulators.

\subsection{Power Dissipation Summary}

\begin{itemize}
    \item \textbf{MOSFET (Q2):} Dissipates $\approx 4.5\text{W}$ at $15\text{A}$ ($25^\circ\text{C}$). Due to $R_{DS(on)}$ increase at high temperature, real dissipation can reach $\approx 7.2\text{W}$. \textbf{A heatsink is mandatory.}
    
    \item \textbf{Zener Diode (D2):} With $R2 = 10\text{k}\Omega$ and $V_{batt} = 24\text{V}$, the current $I_Z$ is only $0.9\text{mA}$. 
    $$P_{D2} = 15\text{V} \cdot 0.9\text{mA} = 13.5\text{mW}$$
    This is far below the $1\text{W}$ rating. Power dissipation is \textbf{negligible}.

    \item \textbf{TVS Diode (D3):} Dissipates \textbf{0W} during normal operation ($V_{batt} < 18\text{V}$). It is rated for $600\text{W}$ peak pulse power to handle millisecond transients.
\end{itemize}

\section{Balancing check}
The PCB is intended to monitor the state of each battery cell and ensure they are balanced.
If a cell voltage exceeds or drops below a certain threshold, a warning will be triggered.

\subsection{Cell 1}
\subsubsection{Specifications}
\begin{itemize}
\item Series voltage: $V_{in} = V_{cell1} = 4.2V (\max) / 3.5V (\min)$
\item $V_{out} = 2.5V (\max)$
\item $I = 100\mu A$ (target current through the voltage divider)
\end{itemize}
The input impedance of the ESP32 ADC is sufficiently high ($>1M\Omega$) to neglect the loading effect on the voltage divider.

\subsubsection{Divider calculation}
Using a voltage divider:
$$V_{out} = V_{in} \cdot \frac{R_2}{R_1 + R_2}$$
$$R_{total} = R_1 + R_2 = \frac{V_{in}}{I} = \frac{4.2V}{100\mu A} = 42k\Omega$$
$$R_2 = \frac{V_{out}}{V_{in}} \cdot R_{total} = \frac{2.5V}{4.2V} \cdot 42k\Omega = 25k\Omega$$
$$R_1 = R_{total} - R_2 = 42k - 25k = 17k\Omega$$

\subsubsection{Standard values (E96 1\%)}
\begin{itemize}
\item $R_1 = 16.9k\Omega$
\item $R_2 = 24.9k\Omega$
\end{itemize}

\subsubsection{Min/Max check}
At $V_{in} = 4.2V$:
$$V_{out} = 4.2V \cdot \frac{24.9k}{16.9k + 24.9k} = 2.50V$$
At $V_{in} = 3.5V$:
$$V_{out} = 3.5V \cdot \frac{24.9k}{16.9k + 24.9k} = 2.08V$$

\subsubsection{Tolerance (1\%)}
At $V_{in} = 4.2V$, worst case scenarios:
\begin{itemize}
    \item Max Output: $R_1 (-1\%), R_2 (+1\%) \rightarrow V_{out} = 2.52V$
    \item Min Output: $R_1 (+1\%), R_2 (-1\%) \rightarrow V_{out} = 2.47V$
\end{itemize}
The output remains safely within the ESP32 ADC linear range.

\subsection{Cell 2}
\subsubsection{Specifications}
\begin{itemize}
\item Series voltage: $V_{in} = V_{cell1} + V_{cell2} = 8.4V (\max) / 7.0V (\min)$
\item $V_{out} = 2.5V (\max)$
\item $I = 100\mu A$ (target current)
\end{itemize}

\subsubsection{Divider calculation}
$$R_{total} = \frac{8.4V}{100\mu A} = 84k\Omega$$
$$R_2 = \frac{2.5V}{8.4V} \cdot 84k\Omega = 25k\Omega$$
$$R_1 = 84k - 25k = 59k\Omega$$

\subsubsection{Standard values (E96 1\%)}
\begin{itemize}
\item $R_1 = 59.0k\Omega$
\item $R_2 = 24.9k\Omega$
\end{itemize}

\subsubsection{Min/Max check}
At $V_{in} = 8.4V$: $V_{out} = 2.49V$ \\
At $V_{in} = 7.0V$: $V_{out} = 2.07V$

\subsubsection{Tolerance check}
Max $V_{out}$ (with $1\%$ tolerance): $2.53V$. Acceptable.

\subsection{Cell 3}
\subsubsection{Specifications}
\begin{itemize}
\item Series voltage: $V_{in} = \sum_{i=1}^{3} V_{celli} = 12.6V (\max) / 10.5V (\min)$
\item $V_{out} = 2.5V (\max)$
\end{itemize}

\subsubsection{Divider calculation}
$$R_{total} = \frac{12.6V}{100\mu A} = 126k\Omega$$
$$R_2 = \frac{2.5V}{12.6V} \cdot 126k\Omega = 25k\Omega$$
$$R_1 = 126k - 25k = 101k\Omega$$

\subsubsection{Standard values (E96 1\%)}
\begin{itemize}
\item $R_1 = 100k\Omega$
\item $R_2 = 24.9k\Omega$
\end{itemize}

\subsubsection{Min/Max check}
At $V_{in} = 12.6V$: $V_{out} = 2.51V$ \\
At $V_{in} = 10.5V$: $V_{out} = 2.09V$

\subsubsection{Tolerance check}
Max $V_{out}$ (with $1\%$ tolerance): $2.55V$. Acceptable.

\subsection{Cell 4}
\subsubsection{Specifications}
\begin{itemize}
\item Series voltage: $V_{in} = \sum_{i=1}^{4} V_{celli} = 16.8V (\max) / 14.0V (\min)$
\item $V_{out} = 2.5V (\max)$
\end{itemize}

\subsubsection{Divider calculation}
$$R_{total} = \frac{16.8V}{100\mu A} = 168k\Omega$$
$$R_2 = \frac{2.5V}{16.8V} \cdot 168k\Omega = 25k\Omega$$
$$R_1 = 168k - 25k = 143k\Omega$$

\subsubsection{Standard values (E96 1\%)}
\begin{itemize}
\item $R_1 = 143k\Omega$
\item $R_2 = 24.9k\Omega$
\end{itemize}

\subsubsection{Min/Max check}
At $V_{in} = 16.8V$: $V_{out} = 2.49V$ \\
At $V_{in} = 14.0V$: $V_{out} = 2.08V$

\subsubsection{Tolerance check}
Max $V_{out}$ (with $1\%$ tolerance): $2.53V$. Acceptable.

\end{document}